% ============================================================
%  五、模型的建立与求解
% ============================================================
\section{模型的建立与求解}

% 这是论文的核心部分，按问题顺序逐一展开。
% 每个问题建议包含：建模思路 → 数学推导 → 算法设计 → 求解过程。

\subsection{问题一：[在此填写问题简称]}

\subsubsection{建模思路}
在此说明本问题的建模动机与核心思想。

\subsubsection{模型建立}
在此给出详细的数学建模过程，包括关键公式。

\subsubsection{模型求解}
在此说明用到的算法、参数设置及求解流程。

% 流程图插入示例：导出 PDF/PNG 后放入 figures/ 目录即可自动替换占位框
\figplot{问题一建模流程图}{4.8cm}{fig:q1-flow}{%
  建议导出为 fig\_q1\_flow.pdf（或 .png），展示数据预处理→建模→求解→验证}{fig_q1_flow.pdf}{%
  节点说明应与正文「模型建立」「模型求解」小节呼应。}

\subsection{问题二：[在此填写问题简称]}

\subsubsection{建模思路}
在此说明本问题的建模动机与核心思想。

\subsubsection{模型建立}
在此给出详细的数学建模过程，包括关键公式。

\subsubsection{模型求解}
在此说明用到的算法、参数设置及求解流程。

% 更多问题可按需复制上述结构
